\phantomsection
\color{unimain}\section*{\centering{\large{ПЕРЕЧЕНЬ СОКРАЩЕНИЙ}}}
\addcontentsline{toc}{section}{Перечень сокращений}
{\color{white}\fontsize{0.1pt}{0.1pt}\selectfont<begABBRS>}
\color{black}
\begin{longtable}{>{\raggedright\arraybackslash}m{2cm}>{\raggedright\arraybackslash}m{0.5cm}>{\raggedright\arraybackslash}m{20cm}}
\endfirsthead\endhead\endfoot\endlastfoot
IEC & -- &  International Electrotechnical Commission; \\
PPS & -- & Pulse Per Second; \\
АПВ & -- & автоматическое повторное включение; \\
АПК & -- & автоматическая проверка канала; \\
АУ & -- & автоматическое ускорение; \\
АУВ & -- & автоматика управления выключателем; \\
БИ & -- & блок испытательный; \\
БК & -- & блокировка при качаниях; \\
БНН & -- & блокировка при неисправности цепей напряжения; \\
ВО & -- & вторичная обмотка; \\
ВЧ & -- & высокочастотный; \\
ДВ & -- & дискретный вход; \\
ДЗ & -- & дистанционная защита; \\
ДФЗ & -- & дифференциально-фазная защита; \\
ЕЭС & -- & единая энергетическая система; \\
ЗНР & -- & защита от неполнофазного режима; \\
ЗНФ & -- & защита от непереключения фаз; \\
ИО & -- & измерительный орган / избирательный орган; \\
ИЧМ & -- & интерфейс человек-машина; \\
ИЭУ & -- & интеллектуальное электронное устройство; \\
КЗ & -- & короткое замыкание; \\
ЛВ & -- & логика включения; \\
ЛС & -- & логика связи; \\
ЛЭП & -- & линия электропередачи; \\
МТЗ & -- & максимальная токовая защита; \\
МФТО & -- & междуфазная токовая отсечка; \\
НКУ & -- & низковольтное комплектное устройство; \\
НПП & -- & научно-производственное предприятие; \\
НТД & -- & нормативно-техническая документация; \\
ОВ & -- & оперативный вывод / обходной выключатель; \\
ОМП & -- & определение места повреждения; \\
ОПУ & -- & общеподстанционный пункт управления; \\
ОСШ & -- & обходная система шин; \\
ОТФ & -- & отключение трех фаз; \\
ОУ & -- & оперативное ускорение; \\
ПК & -- & предохранительный клапан; \\
ПО & -- & пусковой орган / программное обеспечение; \\
ПП & -- & прямая последовательность; \\
ПРД & -- & передатчик; \\
ПРМ & -- & приемник; \\
ПС & -- & предупредительная сигнализация; \\
ПУ & -- & пульт управления / поперечное ускорение; \\
ПУЭ & -- & правила устройства электроустановок; \\
РАС & -- & регистратор аварийных событий; \\
РЗ & -- & релейная защита; \\
РЗА & -- & релейная защита и автоматика; \\
РПВ & -- & реле положения «включено»; \\
РПО & -- & реле положения отключено; \\
РУ & -- & распределительное устройство; \\
РФ & -- & реле фиксации; \\
РЭ & -- & руководство по эксплуатации; \\
СИ & -- & средства измерений; \\
СШ & -- & система (секция) шин; \\
ТЗНП & -- & токовая защита нулевой последовательности; \\
ТК & -- & токовый контроль / текущий контроль; \\
ТН & -- & трансформатор напряжения; \\
ТНЗНП & -- & токовая направленная защита нулевой последовательности; \\
ТО & -- & токовая отсечка; \\
ТУ & -- & телеуправление / телеускорение / технические условия; \\
УРОВ & -- & устройство резервирования при отказе выключателя; \\
ЦН & -- & цепи напряжения; \\
ШОН & -- & шкаф отбора напряжения; \\
ШСВ & -- & шиносоединительный выключатель; \\
ШЭТ & -- & шкаф электротехнический типовой; \\
ЭМО & -- & электромагнит отключения. \\
{\color{white}\fontsize{0.1pt}{0.1pt}\selectfont<endABBRS>}
\end{longtable}
